\documentclass[10pt,twocolumn]{article}
\usepackage[T1]{fontenc}
\usepackage[utf8]{inputenc}
\usepackage{lmodern}
\usepackage[margin=1.8cm,top=2.4cm,bottom=2cm,columnsep=0.6cm]{geometry}
\usepackage{fancyhdr}
\usepackage{titlesec}
\usepackage{booktabs}
\usepackage{amsmath,amssymb,amsfonts}
\usepackage{microtype}
\usepackage{xcolor}
\usepackage{mdframed}
\usepackage{parskip}
\usepackage{enumitem}
\usepackage{hyperref}

\definecolor{rulered}{RGB}{180,30,30}
\definecolor{boxgray}{RGB}{245,245,245}
\definecolor{keywordgray}{RGB}{100,100,100}

\pagestyle{fancy}
\fancyhf{}
\fancyhead[L]{\textsc{\small Claude Mag}}
\fancyhead[C]{\small\textit{Machine Learning Research Digest}}
\fancyhead[R]{\small 2026-07-06}
\fancyfoot[C]{\small\thepage}
\renewcommand{\headrulewidth}{0.8pt}

\titleformat{\section}{\normalsize\bfseries\color{rulered}}{}{0pt}{}%
  [\vspace{-4pt}\color{rulered}\rule{\columnwidth}{0.4pt}]
\titlespacing{\section}{0pt}{8pt}{4pt}

\hypersetup{colorlinks=true,urlcolor=rulered,linkcolor=black}

\begin{document}
\setlength{\parindent}{0pt}
\setlength{\parskip}{4pt}

{\small\color{keywordgray}\textbf{Keywords:}
  robotic manipulation \textbullet{}
  vision-language-action models \textbullet{}
  memory benchmarking \textbullet{}
  long-horizon tasks}\par\vspace{4pt}

{\LARGE\bfseries\raggedright
  RoboMME: Benchmarking and Understanding Memory for
  Robotic Generalist Policies}\par\vspace{4pt}

{\large\itshape\raggedright
  16 Tasks Across Four Memory Categories Reveal That Current VLA Memory
  Mechanisms Are Far From Solving Long-Horizon Manipulation}\par\vspace{6pt}

{\small\textbf{By} Yinpei Dai, Hongze Fu, Jayjun Lee, Yuejiang Liu,
  Haoran Zhang, Jianing Yang, Chelsea Finn, Nima Fazeli \& Joyce Chai}\par\vspace{2pt}

{\small\color{keywordgray}arXiv:2603.04639 \textbullet{}
  ICML 2026 Spotlight Presentation}
\par\vspace{2pt}\noindent{\color{rulered}\rule{\columnwidth}{0.6pt}}\vspace{6pt}

Vision-language-action models have made remarkable progress on short-horizon
manipulation, but long-horizon tasks demanding memory remain poorly understood.
RoboMME introduces the first large-scale standardized benchmark for memory in
robotic generalist policies, comprising 16 tasks across temporal, spatial,
object, and procedural memory categories, and a systematic study of 14
memory-augmented VLA variants. The result: every variant fails significantly on
tasks requiring 5+ step sequences, with the oracle gap exceeding 45\%.

\begin{mdframed}[backgroundcolor=boxgray,linecolor=rulered,linewidth=1pt,
    innerleftmargin=6pt,innerrightmargin=6pt,
    innertopmargin=6pt,innerbottommargin=6pt]
\textbf{\small AT A GLANCE}\par\vspace{4pt}
\begin{description}[leftmargin=0pt,labelsep=4pt,itemsep=2pt,parsep=0pt,
    font=\small\bfseries]
  \item[Problem:] VLA evaluations focus on memoryless tasks; no standardized
    benchmark assesses how well different memory architectures support
    long-horizon, history-dependent manipulation.
  \item[Why it matters:] Real-world manipulation involves counting actions,
    retrieving occluded objects, and executing multi-step procedures ---
    all fundamentally memory-dependent.
  \item[Method:] 16 tasks classified into four memory types; 14 VLA variants
    built on the $\pi_{0.5}$ backbone systematically covering representation
    and integration strategy choices.
  \item[Result:] FrameSamp+Modul achieves best aggregate performance;
    all variants fail on 5+ step procedural sequences.
  \item[Evidence:] Evaluated on LIBERO and LIBERO-Plus with procedurally
    generated variants ensuring evaluation diversity.
\end{description}
\end{mdframed}

\section{The Big Picture}

VLA models extend multimodal LLMs with action outputs, achieving impressive
generalization across robot platforms and tasks. Yet every established
benchmark --- LIBERO, SimplerEnv, MetaWorld --- evaluates performance on tasks
solvable from the current observation alone, without reference to history. This
memoryless evaluation regime is convenient but incomplete.

Real-world manipulation is deeply history-dependent. A robot sorting items
must remember which bin each object was assigned to. A robot following a recipe
must track which ingredients have already been added. A robot retrieving objects
from a drawer must remember where the object was before it was occluded. Without
a principled benchmark for these capabilities, the field has no way to measure
progress on long-horizon robustness.

RoboMME fills this gap.

\section{Why Existing Approaches Fall Short}

\textbf{Benchmark limitations:} Existing benchmarks either test single
tasks independently (no temporal dependency) or use cumulative success rate
on fixed-length sequences, conflating task difficulty with memory requirement.
RoboMME's taxonomy separates these confounds.

\textbf{Architecture comparisons:} Prior work comparing memory mechanisms
(frame stacking, recurrent states, external memory) uses different tasks,
backbones, and training procedures, making conclusions non-transferable.
RoboMME fixes the backbone ($\pi_{0.5}$) and training recipe, varying only
the memory component.

\section{Core Mechanism}

\textbf{Memory taxonomy (four types):}
\begin{itemize}[itemsep=2pt,parsep=0pt]
  \item \textbf{Temporal memory:} Track a sequence of subgoals over time
    (e.g., place five objects in a specified order; count repeated actions).
  \item \textbf{Spatial memory:} Remember and retrieve the location of objects
    that become temporarily occluded during the episode.
  \item \textbf{Object memory:} Identify and track specific objects across
    visual changes (different lighting, partial occlusion, displacement).
  \item \textbf{Procedural memory:} Execute a multi-step procedure from
    an instruction, where the full procedure is not visible in the
    current observation.
\end{itemize}

\textbf{Memory mechanism variants (14):} Each variant combines a
\emph{representation} (FrameSamp with $k \in \{1,4,8\}$ frames; recurrent
hidden state; explicit key-value external memory) with an \emph{integration
strategy} (early concatenation; modulation of the VLM trunk; late fusion into
the action expert). The full factorial design yields 14 meaningful combinations.

\section{Technical Formulation}

\textbf{Memory score:} For task $\tau$ and memory category $c$,
\[
S_\tau = \frac{1}{N_\tau}\sum_{i=1}^{N_\tau} \mathbf{1}[\text{episode } i \text{ succeeds}], \quad
M_c = \frac{1}{|T_c|}\sum_{\tau \in T_c} S_\tau.
\]

\textbf{Memory dependence gap:} For each task $\tau$,
\[
\Delta_\tau = S_\tau(\text{memory agent}) - S_\tau(\text{Markov agent}),
\]
where the Markov agent uses only the current observation (history length 1).
Tasks with $\Delta_\tau > 0.1$ are classified as genuinely memory-dependent.

\textbf{Oracle gap:} The gap to an oracle with perfect recall of all history,
\[
\Delta^\text{oracle}_\tau = S_\tau(\text{oracle}) - \max_\text{variant} S_\tau(\text{variant}).
\]
For procedural memory tasks, $\Delta^\text{oracle} \approx 0.45$ on average ---
current mechanisms capture at most 55\% of the available memory signal.

\textbf{FrameSamp+Modul} cross-attends $k$ sampled historical frames into the
VLM backbone via:
\[
\mathbf{h}_t = \text{VLM}(o_t) + \text{CrossAttn}(\text{VLM}(o_t),\,\{h_{t_j}\}_{j=1}^k),
\]
where $\{t_j\}$ are selected historical timesteps.

\section{Learning or Inference Procedure}

All 14 variants share the same training recipe:
\begin{enumerate}[itemsep=2pt,parsep=0pt]
  \item Pre-train on heterogeneous robot demonstration data (OXE-Mix).
  \item Fine-tune on RoboMME task demonstrations with per-variant
    memory module.
  \item Evaluate on held-out procedurally generated episodes (100 per task).
  \item Report $S_\tau$ and $\Delta_\tau$ for all tasks; $M_c$ per category.
\end{enumerate}

\section{What the Guarantee Says}

RoboMME does not provide formal guarantees, but establishes two empirical
claims with high statistical confidence ($p < 0.01$ across all reported gaps):
(1) 12 of 16 tasks show $\Delta_\tau > 0.1$, confirming genuine memory
dependence. (2) No variant achieves $S_\tau > 0.6$ on tasks requiring 5+ step
procedural sequences.

\section{Experimental Findings}

\begin{itemize}[itemsep=2pt,parsep=0pt]
  \item \textbf{FrameSamp+Modul:} Best aggregate $M_c$ across all four
    categories; particularly strong on temporal counting and time-sensitive tasks.
  \item \textbf{Recurrent:} Best on procedural memory tasks (leverages
    hidden state compression), but degrades on temporal counting.
  \item \textbf{ExternalMem:} Marginal gains over FrameSamp on most tasks;
    naive key-value memory is not yet a silver bullet.
  \item \textbf{Dense history (FrameSamp $k=8$):} Surprisingly weak;
    long raw history contexts confuse the VLM backbone without selective
    attention.
  \item \textbf{All variants:} Fail significantly on 5+ step procedural
    sequences, with oracle gap $>$45\%.
\end{itemize}

\section{Ablations and Interpretation}

\textbf{History length $k$:} Performance peaks at $k=4$ for most tasks;
$k=8$ hurts due to context confusion. Task-adaptive $k$ selection (based
on estimated episode length) is a promising direction.

\textbf{Integration strategy:} Modulation consistently outperforms early
concatenation and late fusion, suggesting that injecting historical context
into intermediate VLM representations is more effective than prepending raw
tokens or appending to the action expert.

\textbf{What memory helps most:} Spatial and object memory benefit most
from frame sampling; procedural memory benefits most from recurrence.
A hybrid design adapting the memory mechanism to the inferred task type
is a natural future direction.

\section*{Reference}

Yinpei Dai, Hongze Fu, Jayjun Lee, Yuejiang Liu, Haoran Zhang, Jianing Yang,
Chelsea Finn, Nima Fazeli, Joyce Chai.
\textbf{RoboMME: Benchmarking and Understanding Memory for
Robotic Generalist Policies}.
arXiv:2603.04639, March 2026.
\url{https://arxiv.org/abs/2603.04639}

\end{document}
